\section{Nutzung der \LaTeX{}-Umgebung}
\label{latex}
Der Fachschaftsrat nutzt \LaTeX{}, um Dokumente aller Art zu verfassen. Grundlage dafür ist das Basispaket \texttt{fsrbase}, auf dem alle Dokumentenklassen des Fachschaftsrats aufbauen. Es sollte immer verwendet werden, wenn ein Dokument mit \LaTeX{} erstellt wird. Darauf aufbauend gibt es Dokumentenklassen, die Layout, Gestaltung und wiederkehrende Elemente eines Dokumenttyps festlegen, sodass sich der Verfasser nur noch um den Inhalt kümmern muss -- beispielsweise bei \hyperref[protokoll]{Sitzungsprotokollen}.

Dieser Abschnitt beschreibt, welche Vorlagen es gibt, was sie festlegen und wie Dokumente auf dem Arbeitsrechner des Fachschaftsrats kompiliert (also die PDF-Dateien erstellt) werden.

\subsection{Übersicht über Paket und Dokumentenklassen}
\label{latex-uebersicht}
\begin{center}
	\begin{tabularx}{\linewidth}{lX}
		\toprule
		\textbf{Name}               & \textbf{Zweck}                                                                                                                   \\
		\midrule
		\texttt{fsrbase}            & \hyperref[latex-fsrbase]{Basispaket}: legt Schriftarten, Farben, deutsche Sprache und Anführungszeichen fest und stellt den Befehl \texttt{\textbackslash fsrlogo} bereit. \\
		\texttt{fsrprotocol}        & Dokumentenklasse für \hyperref[protokoll]{Sitzungsprotokolle}.                                                   \\
		\bottomrule
	\end{tabularx}
\end{center}
% TODO: Weitere Dokumentenklassen ergänzen, sobald diese existieren.

Welche Dokumentenklasse verwendet wird, legt der Befehl \mintinline{latex}{\documentclass} in der ersten Zeile des Dokuments fest. Ein Sitzungsprotokoll beginnt beispielsweise mit \mintinline{latex}{\documentclass{fsrprotocol}}. Das Basispaket \texttt{fsrbase} wird dabei bei allen Dokumentenklassen des Fachschaftsrats automatisch mitgeladen.

\subsection{Sonderzeichen}
\label{latex-sonderzeichen}
Einige Zeichen haben in \LaTeX{} eine besondere Bedeutung, z.\,B. leitet \texttt{\%} einen Kommentar ein. Sie können deshalb nicht einfach so im Text verwendet werden. Sollen sie im Text erscheinen, muss ein Backslash davor stehen, z.\,B. \mintinline{latex}{50\%} oder \mintinline{latex}{Q\&A}.

\begin{center}
	\begin{tabularx}{\linewidth}{clX}
		\toprule
		\textbf{Zeichen} & \textbf{Eingabe}            & \textbf{Bedeutung ohne Backslash}        \\
		\midrule
		\texttt{\%}      & \texttt{\textbackslash\%}   & Leitet einen Kommentar ein               \\
		\texttt{\&}      & \texttt{\textbackslash\&}   & Trennt Spalten in Tabellen               \\
		\texttt{\_}      & \texttt{\textbackslash\_}   & Tiefstellung in Formeln                  \\
		\texttt{\#}      & \texttt{\textbackslash\#}   & Platzhalter für Argumente von Befehlen   \\
		\texttt{\$}      & \texttt{\textbackslash\$}   & Beginnt und beendet eine Formel          \\
		\bottomrule
	\end{tabularx}
\end{center}

\subsection{Das Basispaket \texttt{fsrbase}}
\label{latex-fsrbase}
Das Basispaket sorgt dafür, dass alle Dokumente des Fachschaftsrats einheitlich aussehen. Es legt Folgendes fest:
\begin{itemize}
	\item \textbf{Kompilierprogramm:} Das Paket setzt LuaLaTeX voraus und bricht mit einer \hyperref[latex-fehlermeldungen]{Fehlermeldung} ab, wenn ein anderes Programm (z.\,B. pdfLaTeX) verwendet wird.
\item \textbf{Schriftarten:} TeX Gyre Termes als Serifenschrift und TeX Gyre Heros als serifenlose Schrift (über \texttt{fontspec}).
	\item \textbf{Sprache:} Deutsche Silbentrennung und Typografie über \texttt{babel} mit der Option \texttt{ngerman}.
	\item \textbf{Anführungszeichen:} Über \texttt{csquotes} werden gerade Anführungszeichen \texttt{\string"...\string"} im Text automatisch zu deutschen Anführungszeichen "..." umgewandelt. Alternativ kann \mintinline{latex}{\enquote{...}} verwendet werden.
	\item \textbf{Farben:} Die gewählten Farben des Fachschaftsrats, \texttt{fsrcolor} und \texttt{fsrgrey}, stehen über \texttt{xcolor} zur Verfügung, z.\,B. mit \mintinline{latex}{\textcolor{fsrgrey}{Text}}.
	\item \textbf{Logo:} Der Befehl \mintinline{latex}{\fsrlogo[Höhe]} setzt das Fachschaftslogo. Ohne Angabe der Höhe wird das Logo 5\,cm hoch gesetzt.
\end{itemize}

Außerdem lädt das Basispaket die Pakete \texttt{microtype} (bessere Typografie/Blocksatz) und \texttt{graphicx} (Einbinden von Bildern), die deshalb nicht erneut eingebunden werden müssen.

\subsection{Erstellen von Dokumenten ohne passende Dokumentenklasse}
\label{latex-ohne-dokumentenklasse}
Für Dokumente, für die es noch keine eigene Dokumentenklasse gibt (z.\,B. Aushänge oder Schreiben), kann und soll das Basispaket auch direkt eingebunden werden.

\notebox{Zur Erstellung von Dokumenten sollten stets die \KOMAScript{}-Dokumentenklassen benutzt werden, um einen einheitlichen Stil beizubehalten. Statt \texttt{article} wird also \texttt{scrartcl} verwendet, statt \texttt{report} \texttt{scrreprt}, statt \texttt{book} \texttt{scrbook} und statt \texttt{letter} \texttt{scrlttr2}.}

\begin{Codeblock}[title={Minimales Dokument mit dem Basispaket}]{latex}
	\documentclass{scrartcl}
	\usepackage{fsrbase}

\begin{document}
\fsrlogo[4cm]

\section*{Überschrift}
Hier steht der Inhalt des Dokuments.
\end{document}
\end{Codeblock}

\subsection{Kompilieren einer PDF-Datei}
\label{latex-kompilieren}
Auf dem Arbeitsrechner des Fachschaftsrats werden Dokumente in VSCodium mit der Erweiterung LaTeX Workshop kompiliert:

\begin{enumerate}
	\item Die \texttt{.tex}-Datei in VSCodium öffnen.
	\item Das Dokument mit \texttt{Strg+Alt+B} oder über den grünen Pfeil \textit{Build LaTeX project} oben rechts kompilieren. Standardmäßig wird das Dokument außerdem bei jedem Speichern automatisch neu kompiliert.
	\item Die erzeugte PDF-Datei mit \texttt{Strg+Alt+V} oder über das Symbol \textit{View LaTeX PDF} oben rechts in einer Vorschau neben dem Quelltext öffnen. Die Vorschau aktualisiert sich nach jedem Kompilieren automatisch.
\end{enumerate}
% TODO: Prüfen, ob Tastenkürzel und automatisches Kompilieren beim Speichern auf tesla so eingerichtet sind.

Die PDF-Datei wird im automatisch erstellten Ordner \texttt{build/} abgelegt. Die fertige PDF-Datei kann anschließend beliebig verschoben werden.

Auf dem Arbeitsrechner ist LaTeX Workshop bereits für LuaLaTeX eingerichtet.

\subsection{Erstellen von Dokumenten auf dem eigenen Rechner}
\label{latex-eigener-rechner}
Dokumente können auch auf einem eigenen Rechner kompiliert werden. Dort muss ebenfalls LuaLaTeX verwendet werden (z.\,B. mit dem Terminal-Befehl \texttt{latexmk -lualatex datei.tex}). Damit \LaTeX{} die Vorlagen des Fachschaftsrats findet, müssen folgende Dateien im selben Ordner wie die \texttt{.tex}-Datei liegen:
\begin{itemize}
	\item das Basispaket \texttt{fsrbase.sty},
	\item die jeweils verwendete Dokumentenklasse, z.\,B. \texttt{fsrprotocol.cls},
	\item das Logo \texttt{fsr-logo.png}, falls der Befehl \mintinline{latex}{\fsrlogo} verwendet wird (die Dokumentenklassen tun dies in der Regel).
\end{itemize}

\notebox{Alternativ können die Dateien auch systemweit installiert werden, sodass sie in jedem Ordner zur Verfügung stehen und nicht in jeden Arbeitsordner kopiert werden müssen. Wie das funktioniert, hängt von der installierten \LaTeX{}-Distribution ab (z.\,B. TeX~Live oder MiKTeX). Details beschreibt die Dokumentation der jeweiligen Distribution.}

Bei Online-Editoren wie Overleaf ist eine systemweite Installation nicht möglich. Dort müssen die genannten Dateien immer in das jeweilige Projekt hochgeladen werden und neben der \texttt{.tex}-Datei liegen. Außerdem muss in den Projekteinstellungen (\textit{Menü} \textrightarrow{} \textit{Compiler}) LuaLaTeX als Compiler ausgewählt werden.

\subsection{Fehlermeldungen der Vorlagen}
\label{latex-fehlermeldungen}
Fehler werden in LaTeX Workshop im Reiter \textit{Probleme} am unteren Bildschirmrand angezeigt. Neben den üblichen \LaTeX{}-Fehlern können die Vorlagen des Fachschaftsrats eigene Fehlermeldungen ausgeben:
\begin{description}
	\item[\texttt{LuaTeX is required to compile this document}] Das Dokument wurde nicht mit LuaLaTeX kompiliert (siehe \autoref{latex-kompilieren} bzw. \autoref{latex-eigener-rechner}).
	\item[\texttt{Makro ... ist nicht definiert}] In einem Sitzungsprotokoll fehlt eine der Pflichtangaben am Anfang des Dokuments (siehe \autoref{protokoll-schreiben}).
\end{description}

Bei anderen Fehlern hilft folgendes Vorgehen:
\begin{enumerate}
	\item Im Reiter \textit{Probleme} auf die Fehlermeldung klicken. Der Editor springt dann zu der Zeile, in der \LaTeX{} den Fehler gefunden hat. Die eigentliche Ursache liegt oft in dieser oder in der Zeile davor, z.\,B. eine vergessene Klammer.
	\item Lässt sich der Fehler so nicht finden, die letzten Änderungen rückgängig machen (\texttt{Strg+Z}), bis das Dokument wieder kompiliert, und die Änderung erneut vornehmen.
	\item Hilft das nicht weiter, kann der aktuelle Administrator weiterhelfen.
\end{enumerate}
